% !TeX root = ../../Book1.tex
% 纲要标题：## 第19章　有限 Galois 理论
% 纲要位置：计划纲要.md，第 644 行。

\chapter{有限 Galois 理论}
\label{b1:ch:19}
\BookChapterNumber{b1:ch:19}

\begin{chapterabstract}
有限 Galois 理论用域自同构群研究扩张中的中间域。本章以正规性和可分性确定适用范围，证明 Artin 不动域定理与有限 Galois 对应，并在具体多项式的分裂域中计算这一对应。
\end{chapterabstract}

\begin{learninggoals}
\begin{itemize}
\item 证明 Artin 不动域定理，说明嵌入独立性和线性关系下降分别承担哪个次数估计。
\item 使用 Galois 对应计算中间域次数，辨认正规子扩张并构造限制同态。
\item 从分裂域、根的置换和判别式识别低次方程的 Galois 群。
\item 正确计算合成域与交域，区别实际子域、共轭子域和抽象同构类型，并用正规化子计算中间域的自同构。
\end{itemize}
\end{learninggoals}

在 \(L=\mathbb Q(\sqrt2,\sqrt3)\) 中，可以分别改变两个平方根的符号而保持有理数不动。这给出四个域自同构。其中固定 \(\sqrt2\) 的两个自同构组成一个子群，其共同不动域为 \(\mathbb Q(\sqrt2)\)；固定 \(\sqrt6\) 的两个自同构则给出共同不动域 \(\mathbb Q(\sqrt6)\)。于是同一扩张中的子域，可以通过“哪些对称使其中每个元素不变”来辨认。子域包含的元素越多，要逐点固定它就须满足越多约束，相应的自同构群便越小。一般有限扩张是否也有足够多的自同构，使所有中间域都能如此找回？正规性、可分性与不动域定理将给出这一对应成立的条件。

\ProseParagraphBreak

不动域定理与 Galois 对应使用\cref{b1:thm:embedding-extension}的嵌入延拓、\cref{b1:thm:normal-embedding-criterion}的正规性判据、\cref{b1:thm:separable-embeddings}的嵌入计数和\cref{b1:cor:embedding-independence}的独立性。具体多项式的群计算还使用\secref{b1:sec:1803}的判别式，以及有限域的 Frobenius、中国剩余定理和\cref{b1:lem:integral-closure-elementary}来证明模素数循环型判据。群论方面需要子群指数、正规子群、商群与置换作用。

\ProseParagraphBreak
Galois 对应及多项式的 Galois 群计算，可参阅 Roman 的《Field Theory》\cite[Sections 6.2--6.5, 7.1, 7.4--7.5]{Roman2006FieldTheory}。Lang 的《Algebra》\cite[Chapter VI, \S\S 1--2]{Lang2002}可供比较域与自同构群的结构关系；阅读具体算例时，宜把群的计算和中间域的实际生成元一并写出。

% !TeX root = ../../Book1.tex
% 纲要标题：### 19.1 Galois 扩张
% 纲要位置：计划纲要.md，第 646 行。

\section{Galois 扩张}
\label{b1:sec:1901}

% 纲要内容（仅作写作计划，尚非教材正文）：
%
% * 正规且可分的有限扩张
% * 自同构群与不动域
% * Artin 不动域定理；展开最少非零坐标、归一化与系数下降，并在映射空间中解释独立性给出的维数估计
% * 调用第18.1节已证的域嵌入线性独立性；先完成有限自同构群 \(G\) 作用下 \([L:L^G]=|G|\) 的证明，再进入 Galois 对应
% * 乘法特征与域嵌入独立性的主要证明统一放在第18.1节；本节专门承担 Artin 不动域定理，说明线性独立性在扩张次数估计中的作用
%

多项式根的置换并不都能保持根之间的代数关系。域自同构挑出的正是那些合法的置换。有限 Galois 理论的第一步，是说明什么时候自同构足够多，以及一组有限自同构能够固定多大的子域。

\subsection{有限正规可分扩张}
\label{b1:subsec:1901:01}

\begin{definition}\label{b1:def:finite-galois-extension}
设 \(L/K\) 为有限域扩张。如果 \(L/K\) 同时可分且正规，即每个元素在 \(K\) 上的最小多项式没有重根，并且每个在 \(L\) 中有根的 \(K\) 上不可约多项式都在 \(L\) 中分裂，则称 \(L/K\) 为\term[Galois-kuozhang]{Galois 扩张}[Galois extension]。
\end{definition}
这两个条件分别处理嵌入的数量与像的位置。可分性使最小多项式的根互不相同，保证到代数闭包的不同 \(K\) 嵌入数达到扩张次数；正规性则保证这些嵌入把整个 \(L\) 映到自身。因此两者结合，才把足够多的共轭根选择变成 \(L\) 的自同构。

\ProseParagraphBreak
可分多项式 \(f\in K[x]\) 的分裂域在 \(K\) 上有限，由\cref{b1:prop:splitting-field-exists}得到；它由 \(f\) 的根生成，这些根的最小多项式都是 \(f\) 的因子，因而均可分。\cref{b1:thm:separable-embeddings}保证整个扩张可分，而\cref{b1:thm:normal-embedding-criterion}保证分裂域正规，故它为有限 Galois 扩张。反过来，设 \(L/K\) 有限、正规且可分，取有限生成元 \(L=K(a_1,\ldots,a_r)\)。每个 \(m_{a_i,K}\) 都可分；将这些首一不可约多项式中相同者只保留一次，所得乘积 \(f\) 仍可分，因为两个不同的首一不可约多项式不能有共同根。正规性又保证它们的全部根都属于 \(L\)。这些根生成的域因此包含于 \(L\)，同时包含各 \(a_i\)，也就包含 \(K(a_1,\ldots,a_r)=L\)。所以 \(L\) 恰为 \(f\) 的分裂域。这证明了有限 Galois 扩张也可以等价地描述为可分多项式的分裂域。

\ProseParagraphBreak
\(\mathbb Q(\sqrt2)/\mathbb Q\) 为 Galois 扩张，而 \(\mathbb Q(\sqrt[3]2)/\mathbb Q\) 可分但不正规，因为后者为实域，未包含 \(x^3-2\) 的两个非实根。正特征下，正规性也不能替代可分性：若 \(u\) 在 \(\mathbb F_p\) 上超越，取 \(L=\mathbb F_p(u)\)、\(K=\mathbb F_p(u^p)\)，则 \([L:K]=p\)，且 \(u\) 的最小多项式在 \(L[X]\) 中分解为
\[
X^p-u^p=(X-u)^p.
\]
它的系数属于 \(K\)，全部根已在 \(L\)，并由这些根生成 \(L\)，所以 \(L\) 是它的分裂域，扩张正规。但这 \(p\) 个根全部重合，\(u\) 不可分，故这个扩张不是 Galois 扩张。

\subsection{自同构群与不动域}
\label{b1:subsec:1901:02}

设 \(L/K\) 为域扩张。保持 \(K\) 中每个元素不变的全部 \(L\) 自同构，在复合下组成群 \(\operatorname{Aut}_K(L)\)。这个群对任意扩张都有定义。如果 \(L/K\) 是 Galois 扩张，则把它记为 \(\operatorname{Gal}(L/K)\)，称为扩张 \(L/K\) 的\term[Galois-qun]{Galois 群}[Galois group]。本书在未满足 Galois 条件时仍使用 \(\operatorname{Aut}_K(L)\) 这一记号；这并不表示该扩张没有自同构群。
\symidx{\(\operatorname{Aut}_K(L)\)：保持基域 \(K\) 的 \(L\) 自同构群}
\symidx{\(\operatorname{Gal}(L/K)\)：Galois 扩张的自同构群}

\ProseParagraphBreak
设 \(H\leq\operatorname{Aut}(L)\) 为域 \(L\) 的自同构群的子群。把子域
\[
L^H=\bigl\{\,a\in L\mid \sigma(a)=a\text{ 对每个 }\sigma\in H\,\bigr\}
\]
称为 \(H\) 在 \(L\) 中的\term[budong-yu]{不动域}[fixed field]。
\symidx{\(L^H\)：子群 \(H\) 的不动域}
自同构保持域的运算，因此共同不动元对加、减、乘以及非零元取逆都封闭，确实组成子域。若 \(H\subseteq\operatorname{Aut}_K(L)\)，则 \(K\subseteq L^H\)。

\begin{proposition}\label{b1:prop:galois-automorphism-count}
若 \(L/K\) 为有限 Galois 扩张，则 \(|\operatorname{Gal}(L/K)|=[L:K]\)。
\end{proposition}
\begin{proof}
由\cref{b1:thm:separable-embeddings}，可分性保证 \(L\) 到代数闭包的 \(K\) 嵌入恰有 \([L:K]\) 个。由\cref{b1:thm:normal-embedding-criterion}，正规性保证每个这样的嵌入把 \(L\) 映到自身；有限维单射的像与 \(L\) 维数相同，所以它是自同构。这些嵌入因而恰为 Galois 群的全部元素。
\end{proof}

\subsection{Artin 不动域定理}
\label{b1:subsec:1901:03}

前面从有限扩张 \(L/K\) 出发研究它的自同构；以下\term[Artin-budongyu-dingli]{Artin 不动域定理}[Artin fixed field theorem]反过来，先给域 \(L\) 的有限自同构群 \(G\)，再取 \(K=L^G\)。这里没有预先假设 \(L/K\) 有限：群 \(G\) 的有限性将迫使这个不动域扩张成为有限 Galois 扩张，其次数恰为 \(|G|\)。
\begin{theorem}[Artin]\label{b1:thm:artin-fixed-field}
设 \(G\) 是域 \(L\) 的一个有限自同构群，\(K=L^G\)。则 \(L/K\) 是有限 Galois 扩张，并且
\[
[L:K]=|G|,\qquad \operatorname{Gal}(L/K)=G.
\]
\end{theorem}
\begin{proof}
记 \(G=\{\sigma_1,\ldots,\sigma_n\}\)。次数等式包含两个方向：我们要用线性方程控制 \([L:K]\leq n\)，再用嵌入独立性得到 \(n\leq[L:K]\)。上界的含义是，任意 \(n+1\) 个元素 \(x_1,\ldots,x_{n+1}\in L\) 都在 \(K\) 上线性相关。考察以 \(L\) 为系数域、以 \(c_1,\ldots,c_{n+1}\in L\) 为未知量的齐次方程组：
\[
\sum_{j=1}^{n+1}\sigma_i(x_j)c_j=0\qquad(1\leq i\leq n).
\]
因为只有 \(n\) 个方程，却有 \(n+1\) 个未知量，存在非零的 \(L\) 解。但这些系数暂时未必在 \(K\) 中，还不能用它得到所需的 \(K\) 线性关系。为了使系数下降到不动域，从全部非零解中取非零坐标数最少者，重新编号，使它的非零坐标恰为前 \(r\) 个。再把整个解除以 \(c_r\)，仍得一个解，并可设 \(c_r=1\)。这个归一化使任何 \(\tau\in G\) 都固定该坐标。

将 \(\tau\) 作用到全部方程，得到
\[
\sum_{j=1}^{n+1}(\tau\sigma_i)(x_j)\tau(c_j)=0
\qquad(1\leq i\leq n).
\]
左乘映射 \(\sigma\mapsto\tau\sigma\) 是 \(G\) 的一个排列，逆映射为左乘 \(\tau^{-1}\)；所以上式仅把原来各方程重新排列，向量 \(\bigl(\tau(c_j)\bigr)_j\) 也是原方程组的解。两解之差在第 \(r\) 个坐标为 \(\tau(1)-1=0\)，在 \(r\) 后也均为零。若差非零，它至多有 \(r-1\) 个非零坐标，违背原解的最小选择。因此
\[
\tau(c_j)=c_j\quad(\tau\in G),\qquad c_j\in L^G=K.
\]
取恒等自同构对应的方程，就有 \(\sum_j x_jc_j=0\)，而 \(c_r=1\) 保证这个 \(K\) 系数关系非平凡。任意 \(n+1\) 个元素都如此，故 \([L:K]\leq n\)，特别地，扩张有限。

记 \(d=[L:K]\)，并把所有 \(K\) 线性映射 \(L\rightarrow L\) 组成的空间记为 \(\mathcal V\)。其加法与 \(L\) 标量乘法都是逐点定义的，特别是
\[
(\lambda f)(x)=\lambda f(x)\qquad(\lambda\in L,\ f\in\mathcal V).
\]
乘上 \(\lambda\) 后仍是 \(K\) 线性映射，因为 \(L\) 的乘法交换。若 \(b_1,\ldots,b_d\) 是 \(L/K\) 的一组基，一个 \(K\) 线性映射 \(f\) 完全由 \(f(b_1),\ldots,f(b_d)\in L\) 决定，而且任意指定这 \(d\) 个像都给出唯一的 \(K\) 线性映射。因此求值映射
\[
\mathcal V\longrightarrow L^d,\qquad
f\longmapsto\bigl(f(b_1),\ldots,f(b_d)\bigr)
\]
是 \(L\) 线性同构，故 \(\dim_L\mathcal V=d\)。每个 \(\sigma_i\) 固定 \(K\)，所以属于 \(\mathcal V\)。\cref{b1:cor:embedding-independence}在此说明：若 \(\sum_i\lambda_i\cdot\sigma_i(x)=0\) 对所有 \(x\in L\) 成立，其中 \(\lambda_i\in L\)，则全部 \(\lambda_i=0\)。因此这 \(n\) 个不同自同构是 \(d\) 维 \(L\) 向量空间中的线性无关向量，得到 \(n\leq d\)。结合上界，\(d=n\)。

对任意 \(a\in L\)，把 \(G\) 轨道中的互异元素记作 \(a_1,\ldots,a_s\)，并构造轨道多项式
\[
P_a(t)=\prod_{j=1}^s(t-a_j).
\]
任何 \(\tau\in G\) 都把轨道映到自身且映满，所以只会排列这些根。\(P_a\) 的系数由它们的初等对称式给出，因而被每个 \(\tau\) 固定，属于 \(L^G=K\)。又因为只将互异轨道元素各取一次，\(P_a\) 没有重根，并已在 \(L[t]\) 中分裂。恒等自同构保证 \(a\) 在轨道中，因此 \(P_a(a)=0\)，\(a\) 的最小多项式整除 \(P_a\)，也就可分并在 \(L\) 中分裂。任意 \(a\) 均如此，说明 \(L/K\) 正规且可分。最后由\cref{b1:prop:galois-automorphism-count}，其全部 \(K\) 自同构只有 \(n\) 个；已经包含的 \(G\) 也有 \(n\) 个元素，所以它就是全群。
\end{proof}

若给定的是一个抽象有限群 \(\Gamma\) 对 \(L\) 的作用，作用同态为 \(\rho:\Gamma\rightarrow\operatorname{Aut}(L)\)，则共同不动元只取决于像群 \(\rho(\Gamma)\)。上面的定理应用于这个实际自同构群，给出
\[
[L:L^{\rho(\Gamma)}]=|\rho(\Gamma)|
=\frac{|\Gamma|}{|\ker\rho|}.
\]
最后一个等号是群同态的第一同构定理。只有作用忠实、\(\ker\rho=1\) 时，才能直接以抽象群 \(\Gamma\) 的阶作为扩张次数；这正是前面群作用通过核的商群给出忠实作用的原则。

\subsection{不动域扩张次数的计算}
\label{b1:subsec:1901:04}

对有限 Galois 扩张 \(L/K\)，令 \(G=\operatorname{Gal}(L/K)\)。\cref{b1:thm:artin-fixed-field}与\cref{b1:prop:galois-automorphism-count}分别给出
\[
[L:L^G]=|G|=[L:K].
\]
因 \(K\subseteq L^G\)，塔式公式迫使 \([L^G:K]=1\)，所以 \(L^G=K\)。这说明基域恰能从全部自同构中恢复。

\ProseParagraphBreak
取特征不为二的域 \(k\) 及在 \(k\) 上超越的元素 \(t\)。虽然 \(k(t)/k\) 不是代数扩张，仍可对 \(k(t)\) 的有限自同构群使用不动域定理。令 \(s\) 为逐点固定 \(k\) 且满足 \(s(t)=-t\) 的 \(k(t)\) 自同构，它生成二阶群 \(G=\langle s\rangle\)。群 \(G\) 的不动域包含 \(k(t^2)\)。元素 \(t\) 满足 \(x^2-t^2=0\)，但不属于 \(k(t^2)\)：若 \(t=f(t^2)/g(t^2)\)，其中 \(f,g\in k[x]\)、\(g\ne0\)，则 \(t\cdot g(t^2)=f(t^2)\) 的两侧分别只含奇次幂和偶次幂，且均不为零，矛盾。因此 \([k(t):k(t^2)]=2\)；与不动域定理给出的 \([k(t):k(t)^G]=2\) 比较，得到 \(k(t)^G=k(t^2)\)。

\subsection{自同构数判据与次数估计}
\label{b1:subsec:1901:05}

\begin{corollary}\label{b1:cor:galois-automorphism-criterion}
设 \(L/K\) 为有限域扩张。则 \(L/K\) 为 Galois 扩张，当且仅当 \(|\operatorname{Aut}_K(L)|=[L:K]\)。
\end{corollary}
\begin{proof}
若 \(L/K\) 为 Galois 扩张，结论由\cref{b1:prop:galois-automorphism-count}给出。反过来，令 \(H=\operatorname{Aut}_K(L)\)。假设 \(|H|=[L:K]\)，则\cref{b1:thm:artin-fixed-field}给出 \([L:L^H]=|H|=[L:K]\)。由于 \(K\subseteq L^H\)，塔式公式说明 \(L^H=K\)。该定理还保证 \(L/L^H\) 正规且可分，故 \(L/K\) 为 Galois 扩张。
\end{proof}

\cref{b1:thm:artin-fixed-field}中的两个次数估计有不同来源。线性方程加最短支集论证，把 \(L\) 系数的关系下降为不动域系数的关系，给出上界；嵌入独立性则保证自同构不可能多于维数，给出下界。轨道多项式至多有 \(|G|\) 个根，能说明每个单独的 \(K(a)/K\) 次数至多为 \(|G|\)。但在尚未证明 \(L\) 由一个元素生成时，这只控制各个单扩张，不能直接控制整个 \(L/K\) 的维数。线性方程法处理的是任意 \(|G|+1\) 个元素之间的关系，因而直接控制整个空间，不需要先知道扩张单生成。

% !TeX root = ../../Book1.tex
% 纲要标题：### 19.2 Galois 对应
% 纲要位置：计划纲要.md，第 654 行。

\section{Galois 对应}
\label{b1:sec:1902}

% 纲要内容（仅作写作计划，尚非教材正文）：
%
% * 中间域与子群的反向对应
% * 次数、指数与群阶
% * 正规子群、正规子扩张与商群
% * 紧接对应定理，以 \(X^3-2\) 的分裂域枚举全部子群、不动域、次数、正规性和包含关系
%

现在固定有限 Galois 扩张 \(L/K\)，记 \(G=\operatorname{Gal}(L/K)\)。给定中间域 \(E\)，考察逐点固定 \(E\) 的自同构所组成的子群 \(\operatorname{Gal}(L/E)\)；给定子群，则取其共同不动元。两个构造反转包含关系；要证明它们互逆，还须使用不动域定理中的精确次数。

\subsection{中间域与子群的反向对应}
\label{b1:subsec:1902:01}

下述对应称为\term[Galois-duiying]{Galois 对应}[Galois correspondence]。
\begin{theorem}\label{b1:thm:finite-galois-correspondence}
有限 Galois 扩张 \(L/K\) 的中间域与 \(G=\operatorname{Gal}(L/K)\) 的子群，通过
\[
E\longmapsto\operatorname{Gal}(L/E),\qquad H\longmapsto L^H
\]
建立互逆的、反转包含关系的一一对应。对任意中间域 \(E\)，扩张 \(L/E\) 都是 Galois 扩张。
\end{theorem}
\begin{proof}
先说明最后一句。每个元素在 \(E\) 上的最小多项式整除它在 \(K\) 上的最小多项式，故可分；任意 \(E\) 嵌入也是 \(K\) 嵌入，所以由\cref{b1:thm:normal-embedding-criterion}，它保持 \(L\)，故 \(L/E\) 正规。于是 \(H=\operatorname{Gal}(L/E)\) 满足 \(|H|=[L:E]\)。不动域定理给出 \([L:L^H]=|H|\)，结合 \(E\subseteq L^H\) 以及塔式公式，得到 \(E=L^H\)。

反向给定任意 \(H\leq G\)，\cref{b1:thm:artin-fixed-field}直接给出 \(\operatorname{Gal}(L/L^H)=H\)。这证明两个构造互逆。若 \(E_1\subseteq E_2\)，固定较大域的条件更强，所以 \(\operatorname{Gal}(L/E_2)\subseteq\operatorname{Gal}(L/E_1)\)；子群越大，不动域越小。故对应反转包含。
\end{proof}
不能把最后一句反过来读成每个 \(E/K\) 都正规。例如三次多项式的非正规根域可以作为其 Galois 分裂域的中间域。

\subsection{次数、指数与群阶}
\label{b1:subsec:1902:02}

对对应的 \(E=L^H\)，不动域定理和塔式公式分别给出
\begin{equation}\label{b1:eq:galois-degree-index}
[L:E]=|H|,\qquad [E:K]=[G:H].
\end{equation}
若 \(E_1\subseteq E_2\) 对应 \(H_1\supseteq H_2\)，则 \([E_2:E_1]=[H_1:H_2]\)。这一等式计算的是群的指数，不要求相应子群正规。

\ProseParagraphBreak
例如 \(L=\mathbb F_{q^6}\)、\(K=\mathbb F_q\) 时，\cref{b1:thm:finite-field-automorphisms}给出 \(G=\langle\sigma\rangle\cong C_6\)。指数为 \(d\) 的子群 \(\langle\sigma^d\rangle\) 对应 \(\mathbb F_{q^d}\)，其中 \(d\mid6\)。这与\cref{b1:thm:finite-field-subfields}的直接计算一致；有限域的子域分类并不依赖本定理，因而这里是对两种方法的比较。

\subsection{正规子群、正规子扩张与商群}
\label{b1:subsec:1902:03}

\begin{theorem}\label{b1:thm:galois-quotient}
设 \(L/K\) 为有限 Galois 扩张，\(G=\operatorname{Gal}(L/K)\)，\(H\leq G\)，并令 \(E=L^H\)。则 \(E/K\) 为 Galois 扩张当且仅当 \(H\) 是 \(G\) 的正规子群。此时限制映射 \(G\rightarrow\operatorname{Gal}(E/K)\) 是满同态，且
\[
\operatorname{Gal}(E/K)\cong G/H.
\]
\end{theorem}
\begin{proof}
对任意 \(\sigma\in G\)，直接检验得到
\[
\sigma(E)=L^{\sigma H\sigma^{-1}}.
\]
事实上，若 \(x\in E\)，则 \(\sigma h\sigma^{-1}\) 固定 \(\sigma(x)\)；反向对 \(\sigma^{-1}\) 作同样检验。若 \(H\) 正规，则每个 \(\sigma\) 都保持 \(E\)。任意 \(E\) 的 \(K\) 嵌入能由\cref{b1:thm:embedding-extension}延拓到 \(L\) 的 \(K\) 嵌入，而 \(L/K\) 正规使此延拓属于 \(G\)。它保持 \(E\)，所以嵌入判据说明 \(E/K\) 正规；\(E/K\) 的可分性来自 \(L/K\)，于是为 Galois 扩张。

若 \(E/K\) 正规，每个 \(\sigma\in G\) 都保持 \(E\)，从而 \(L^{\sigma H\sigma^{-1}}=E=L^H\)。Galois 对应的单射性给出 \(\sigma H\sigma^{-1}=H\)，所以 \(H\) 正规。此时限制映射 \(G\rightarrow\operatorname{Gal}(E/K)\) 的核是恰好固定 \(E\) 的 \(H\)；任意目标自同构均能延拓到 \(L\)，故映射满射。\cref{b1:thm:first-isomorphism}遂把商群 \(G/H\) 同构到这个实际像。
\end{proof}
正规性确保限制后仍是 \(E\) 到自身的映射；若 \(E/K\) 不正规，某些 \(\sigma\) 会把它送到另一个共轭子域，不能对整个 \(G\) 定义到 \(\operatorname{Aut}_K(E)\) 的限制同态。

\ProseParagraphBreak
对上述 \(E=L^H\)，逐点固定 \(E\) 与将 \(E\) 保持为集合是不同的条件。记 \(N_G(H)=\{\,\sigma\in G\mid\sigma H\sigma^{-1}=H\,\}\) 为 \(H\) 在 \(G\) 中的正规化子。Galois 对应及证明中的共轭公式分别给出
\[
\{\,\sigma\in G\mid\sigma|_E=\operatorname{id}_E\,\}
=\operatorname{Gal}(L/E)=H,
\qquad
\{\,\sigma\in G\mid\sigma(E)=E\,\}=N_G(H).
\]
恰是 \(N_G(H)\) 中的自同构才能限制成 \(E\) 的 \(K\) 自同构；\cref{b1:ex:19-normalizer-automorphisms}将用它计算 \(\operatorname{Aut}_K(E)\)。

\subsection{中间域的枚举}
\label{b1:subsec:1904:01}

先以 \(X^3-2\) 的分裂域把对应定理用在一个非交换 Galois 群上：求出全部子群及其不动域，并核对次数、正规性和包含关系。

\begin{proposition}\label{b1:prop:cubic-splitting-field-correspondence}
令 \(a=\sqrt[3]2\in\mathbb R\)，\(\zeta\) 为本原三次单位根，并置 \(L=\mathbb Q(a,\zeta)\)。则 \([L:\mathbb Q]=6\)，且 \(G=\operatorname{Gal}(L/\mathbb Q)\cong S_3\)。除 \(\mathbb Q\) 与 \(L\) 外，全部中间域为唯一的二次域 \(\mathbb Q(\zeta)\) 和三个互异的三次域 \(\mathbb Q(\zeta^j a)\)、\(j=0,1,2\)。前者在 \(\mathbb Q\) 上正规，后三者均不正规。
\end{proposition}
\begin{proof}
由\cref{b1:thm:eisenstein}，\(X^3-2\) 在 \(\mathbb Q\) 上不可约，故 \([\mathbb Q(a):\mathbb Q]=3\)。域 \(\mathbb Q(a)\) 包含于 \(\mathbb R\)，而 \(\zeta\) 非实且满足 \(X^2+X+1=0\)，所以 \([L:\mathbb Q(a)]=2\)。这是一个特征零多项式的分裂域，因而为六次 Galois 扩张。群 \(G\) 忠实置换三个根 \(a,\zeta a,\zeta^2a\)，由阶数知它给出全部 \(S_3\)。

记循环置换三个根的自同构为 \(s\)，复共轭为 \(t\)，则
\[
s(a)=\zeta a,\quad s(\zeta)=\zeta,\qquad
t(a)=a,\quad t(\zeta)=\zeta^{-1}.
\]
它们满足 \(s^3=t^2=1\)、\(tst=s^{-1}\)，且生成 \(G\)。\(G\) 的全部子群为 \(G\)、\(\langle s\rangle\)、\(\langle s^{2j}t\rangle\)（\(j=0,1,2\)）及 \(1\)：真非平凡子群只能有二阶或三阶，分别由三个换位或一个三循环生成。

自同构 \(s\) 固定 \(\zeta\)，而 \(s^{2j}t\) 固定 \(\zeta^j a\)。由\eqref{b1:eq:galois-degree-index}，这两个子群的不动域在 \(\mathbb Q\) 上分别有次数二与三；\(\mathbb Q(\zeta)\) 与 \(\mathbb Q(\zeta^j a)\) 也分别有这些次数，故相应包含必为相等。对应的三个二阶子群互异，所以三个三次域也互异。子群 \(\langle s\rangle\) 正规，三个二阶子群彼此共轭而不正规，\cref{b1:thm:galois-quotient}遂给出各中间域的正规性。
\end{proof}

例如取 \(E=\mathbb Q(\zeta)\)。全部六个 \(G\) 中的自同构都将 \(E\) 保持为集合，但只有 \(\langle s\rangle\) 的三个元素逐点固定 \(E\)；复共轭 \(t\) 把 \(\zeta\) 送到 \(\zeta^{-1}\)，所以不逐点固定它。

\cref{b1:tab:cubic-galois-correspondence}集中列出上述对应。对每一行都有 \([E:\mathbb Q]=[G:H]\)、\([L:E]=|H|\)；最后一列同时说明 \(H\) 是否为 \(G\) 的正规子群，因为这里的全部扩张均可分。
\begin{table}[htbp]
\centering
\caption[\(X^3-2\) 分裂域的 Galois 对应]{\(X^3-2\) 分裂域的 Galois 对应。其中 \(j=0,1,2\)。}
\label{b1:tab:cubic-galois-correspondence}
\begin{tabular}{@{}llccl@{}}
\toprule
子群 \(H\) & 不动域 \(E=L^H\) & \([E:\mathbb Q]\) & \([L:E]\) & \(E/\mathbb Q\) 正规 \\
\midrule
\(G\cong S_3\) & \(\mathbb Q\) & \(1\) & \(6\) & 是 \\
\(\langle s\rangle\cong C_3\) & \(\mathbb Q(\zeta)\) & \(2\) & \(3\) & 是 \\
\(\langle s^{2j}t\rangle\cong C_2\) & \(\mathbb Q(\zeta^j a)\) & \(3\) & \(2\) & 否 \\
\(1\) & \(L\) & \(6\) & \(1\) & 是 \\
\bottomrule
\end{tabular}
\end{table}

\cref{b1:fig:cubic-intermediate-fields}将表中的三个三次域分别画出；图中较低的域包含于较高的域，边上的数字是相邻两域的扩张次数。
\begin{figure}[htbp]
\centering
\begin{tikzpicture}[x=1cm,y=1cm,every node/.style={font=\small}]
\node (L) at (0,1.8) {\(L\)};
\node (Z) at (-5.5,0) {\(\mathbb Q(\zeta)\)};
\node (A) at (-1.8,0) {\(\mathbb Q(a)\)};
\node (B) at (1.8,0) {\(\mathbb Q(\zeta a)\)};
\node (C) at (5.5,0) {\(\mathbb Q(\zeta^2a)\)};
\node (Q) at (0,-1.8) {\(\mathbb Q\)};
\draw (L) -- node[pos=.55,fill=white,inner sep=1pt] {\scriptsize 3} (Z);
\draw (L) -- node[pos=.55,fill=white,inner sep=1pt] {\scriptsize 2} (A);
\draw (L) -- node[pos=.55,fill=white,inner sep=1pt] {\scriptsize 2} (B);
\draw (L) -- node[pos=.55,fill=white,inner sep=1pt] {\scriptsize 2} (C);
\draw (Z) -- node[pos=.45,fill=white,inner sep=1pt] {\scriptsize 2} (Q);
\draw (A) -- node[pos=.45,fill=white,inner sep=1pt] {\scriptsize 3} (Q);
\draw (B) -- node[pos=.45,fill=white,inner sep=1pt] {\scriptsize 3} (Q);
\draw (C) -- node[pos=.45,fill=white,inner sep=1pt] {\scriptsize 3} (Q);
\end{tikzpicture}
\caption[\(X^3-2\) 分裂域的中间域格]{\(X^3-2\) 分裂域的中间域格。边上的数字为较高的域在较低的域上的次数。}
\label{b1:fig:cubic-intermediate-fields}
\end{figure}

图中的四个非平凡真中间域两两互不包含：二次与三次之间由次数的整除性排除包含，三个不同的三次域之间也不能包含。把每条包含关系反向，就得到 \(G\) 的子群格。例如 \(\mathbb Q\subset\mathbb Q(\zeta)\subset L\) 对应 \(G\supset\langle s\rangle\supset1\)。这个枚举穷尽所有中间域，因为已经列出了全部子群；一般地，有限 Galois 扩张的中间域数必有限。

% !TeX root = ../../Book1.tex
% 纲要标题：### 19.3 Galois 群的计算与识别
% 纲要位置：计划纲要.md，第 660 行。

\section{Galois 群的计算与识别}
\label{b1:sec:1903}

% 纲要内容（仅作写作计划，尚非教材正文）：
%
% * 二次、多二次扩张；三次式 \(X^3-2\) 的完整对应计算见第19.2节
% * 四次方程的典型 Galois 群
% * 从置换作用和分解型识别子群
% * 模素数分解与循环型：在此完整证明约化判据，第21.3节直接引用
%

识别 Galois 群通常需要两类信息：扩张次数确定群阶，群对根的作用将它表示为对称群的子群。以下每个例子都先明确分裂域，再计算自同构；不能只写出一个抽象群名而略去它怎样作用。

\subsection{二次与多二次扩张}
\label{b1:subsec:1903:01}

特征不为二时，若 \(d\in K^\times\) 不是平方，\(K(\sqrt d)/K\) 为二次 Galois 扩张，非平凡自同构把 \(\sqrt d\) 送到 \(-\sqrt d\)。更一般的可分二次多项式 \(x^2+bx+c\) 的两根之和为 \(-b\)，所以添加一根后另一根已在域内；即使特征为二，可分二次扩张也仍为 Galois 扩张。

\ProseParagraphBreak
考虑 \(L=\mathbb Q(\sqrt2,\sqrt3)\)。若 \(\sqrt3=a+b\sqrt2\)，其中 \(a,b\in\mathbb Q\)，平方比较得 \(2ab=0\)、\(a^2+2b^2=3\)。若 \(b=0\)，三为有理数平方；若 \(a=0\)，\(3/2\) 为有理数平方；两者都与素因数指数的奇偶性矛盾。因此 \([L:\mathbb Q]=4\)，基可取 \(1,\sqrt2,\sqrt3,\sqrt6\)。独立改变两个平方根的符号保持全部乘法关系，给出四个自同构，所以
\[
\operatorname{Gal}(L/\mathbb Q)\cong C_2\times C_2.
\]
三个二元子群的不动域分别为 \(\mathbb Q(\sqrt2)\)、\(\mathbb Q(\sqrt3)\)、\(\mathbb Q(\sqrt6)\)。例如同时改变两者符号的自同构固定 \(\sqrt6\)，而不动域次数为二，故恰为最后一个域。

\ProseParagraphBreak
三次式 \(X^3-2\) 的六次分裂域、\(S_3\) 作用和三个非正规三次根域，已经在\cref{b1:prop:cubic-splitting-field-correspondence}中逐一算出。该例说明，求得 Galois 群以后，还须确定各自同构对根的具体作用，才能识别它固定的域。后文将用判别式区分可分不可约三次式的 \(S_3\) 与 \(C_3\) 情形。

\subsection{四次方程的 Galois 群}
\label{b1:subsec:1903:02}

四次不可约多项式的分裂域次数未必为四。三个典型计算分别给出不同的群。

\ProseParagraphBreak
先取 \(f=x^4-10x^2+1\)。它的根为 \(\pm\sqrt2\pm\sqrt3\)。令 \(a=\sqrt2+\sqrt3\)，则 \(a^{-1}=\sqrt3-\sqrt2\)，所以
\[
\sqrt3=\frac{a+a^{-1}}2,\qquad \sqrt2=\frac{a-a^{-1}}2.
\]
由上一小节的次数计算，\([\mathbb Q(a):\mathbb Q]=4\)。因此 \(f\) 是 \(a\) 的最小多项式，且其分裂域就是 \(\mathbb Q(\sqrt2,\sqrt3)\)，Galois 群为四元群 \(C_2\times C_2\)。

\ProseParagraphBreak
再取 \(f=x^4+ x^3+x^2+x+1\)，其根为非一的五次单位根。变元平移后的 \(f(x+1)\) 由素数五满足\cref{b1:thm:eisenstein}，故 \(f\) 不可约。若 \(\zeta\) 是一个根，所有根 \(\zeta,\zeta^2,\zeta^3,\zeta^4\) 都在 \(\mathbb Q(\zeta)\) 中，扩张次数为四。自同构由 \(\zeta\mapsto\zeta^j\) 给出，复合对应指数模五相乘。模五的二依次给出 \(2,4,3,1\)，所以这个 Galois 群为 \(C_4\)。

\ProseParagraphBreak
最后取 \(f=x^4-2\)，令 \(a=\sqrt[4]2>0\)。分裂域为 \(L=\mathbb Q(a,i)\)。Eisenstein 判据给出实子域 \(\mathbb Q(a)\) 的次数为四；加入 \(i\) 后次数变为八。设 \(r\) 固定 \(i\) 而把 \(a\) 送到 \(ia\)，设 \(s\) 为复共轭。前者确实存在，因为 \([L:\mathbb Q(i)]=4\)，\(a\) 在 \(\mathbb Q(i)\) 上的最小多项式仍是 \(x^4-2\)，可以把 \(a\) 映到另一个根 \(ia\)。直接作用于生成元可验
\[
r^4=s^2=1,\qquad srs=r^{-1}.
\]
四个 \(r^j\) 固定 \(i\)，四个 \(sr^j\) 把 \(i\) 变号，彼此不同，已经列出全部八个自同构。因此群为本书记号下的 \(D_8\)。

\subsection{置换作用、分解型与子群识别}
\label{b1:subsec:1903:03}

若首一可分多项式 \(f\in K[x]\) 的分裂域为 \(L\)，其 Galois 群在各根上的作用忠实，因为固定全部根就固定它们生成的域。一个不可约因子的根构成一个轨道：任意两根由简单扩张同构对应，该同构能由\cref{b1:thm:embedding-extension}延拓到 \(L\)，正规性保证延拓为自同构。因此在 \(K\) 上各不可约因子的次数，正是全部根分成的轨道大小。特别地，不可约多项式给出传递作用。

\ProseParagraphBreak
还可利用判别式辨认置换的奇偶性。
\begin{proposition}[判别式判据]\label{b1:prop:discriminant-alternating-group}
设 \(K\) 是特征不为二的域，\(f\in K[X]\) 是次数为 \(n\geq1\) 的首一可分多项式，\(L\) 是 \(f\) 在 \(K\) 上的分裂域。令 \(G=\operatorname{Gal}(L/K)\)，并通过它在 \(f\) 的全部根上的作用将其视为 \(S_n\) 的子群。则
\begin{equation}\label{b1:eq:discriminant-alternating-group}
G\subseteq A_n\quad\Longleftrightarrow\quad
\operatorname{disc}(f)\text{ 是 }K\text{ 中的平方}.
\end{equation}
\end{proposition}
\begin{proof}
把 \(f\) 的互异根记为 \(a_1,\ldots,a_n\)，并置 \(\delta=\prod_{i<j}(a_j-a_i)\)。交换两个根使 \(\delta\) 变号，因此任意 \(\sigma\in G\) 满足
\[
\sigma(\delta)=\operatorname{sgn}(\sigma)\delta.
\]
\(\delta^2=\operatorname{disc}(f)\) 属于 \(K\)，且它是 \(K\) 中的平方当且仅当 \(\delta\in K\)：若 \(\delta^2=c^2\)，则 \(\delta=\pm c\)；反向显然。由于 \(\delta\ne0\) 且特征不为二，上式又说明 \(G\subseteq A_n\) 当且仅当 \(G\) 的每个元素都固定 \(\delta\)。由 \(L^G=K\) 即得结论。
\end{proof}

特征二的限制确有必要。例如 \(f=X^2+X+1\in\mathbb F_2[X]\) 不可约且可分，其分裂域为 \(\mathbb F_4\)；非平凡自同构交换两个根，故 \(G=S_2\not\subseteq A_2\)。但 \(\operatorname{disc}(f)=1\) 是 \(\mathbb F_2\) 中的平方。此时 \(1=-1\)，从 \(\sigma(\delta)=\operatorname{sgn}(\sigma)\delta\) 不能辨出置换的奇偶性。

若 \(\operatorname{char}K\ne2\)，一个 \(K\) 上的可分不可约三次多项式的 Galois 群是 \(S_3\) 的传递子群，其阶被三整除且整除六，所以只有 \(A_3\cong C_3\) 与 \(S_3\) 两种可能；\cref{b1:prop:discriminant-alternating-group}用判别式是否为平方将它们区分。

\ProseParagraphBreak
这里的轨道大小取自原基域上的不可约分解。若改在有限域中分解整数多项式，则可以借助 Frobenius 自同构得到原 Galois 群中单个元素的循环型。下面建立两者的联系。

\subsection{模素数分解与循环型}
\label{b1:subsec:1903:04}

对一个整数多项式，模素数后的因子次数往往容易计算。要把这些数据用于特征零中的 Galois 群，须使有限域中置换根的 Frobenius 自同构来自原分裂域的自同构。以下证明用全部根生成的整环及其有限剩余域实现这一点；用到的工具是\cref{b1:lem:integral-closure-elementary}的整性封闭性、\cref{b1:thm:crt}的中国剩余定理，以及已经建立的有限域和有限 Galois 理论。

\begin{lemma}\label{b1:lem:reduction-cycle-type}
设 \(f\in\mathbb Z[X]\) 为首一多项式，并且 \(f\) 在 \(\mathbb Q\) 上无重根。令 \(L\) 为 \(f\) 在 \(\mathbb Q\) 上的分裂域，\(G=\operatorname{Gal}(L/\mathbb Q)\)。设素数 \(p\) 使 \(f\) 的模 \(p\) 约化 \(\bar f\in\mathbb F_p[X]\) 也无重根。若 \(\bar f\) 的不同首一不可约因子的次数为 \(d_1,\ldots,d_r\)，则 \(G\) 在 \(f\) 的全部根上的置换作用中包含一个循环长度为 \(d_1,\ldots,d_r\) 的元素。
\end{lemma}
\begin{proof}
设分裂域为 \(L\)，根为 \(\alpha_1,\ldots,\alpha_n\)，\(G=\operatorname{Gal}(L/\mathbb Q)\)，并取环 \(A=\mathbb Z[\alpha_1,\ldots,\alpha_n]\)。所有根都满足首一整数方程，\cref{b1:lem:integral-closure-elementary}说明 \(A\) 由有限多个元素作为 \(\mathbb Z\) 模生成，且每个元素都在 \(\mathbb Z\) 上整。\(G\) 置换这些根，所以保持 \(A\)。由 \(L^G=\mathbb Q\)，\(A^G\) 中的元素为有理整元素，只能是整数：一个既约分数若满足首一整数方程，清除分母后就要求分母整除分子的某次幂，故分母为一。因此 \(A^G=\mathbb Z\)。

理想 \(pA\) 为真理想，否则 \(1/p\in A\)，与刚才的有理整元素判别矛盾。由\cref{b1:prop:maximal-ideal-exists}取极大理想 \(\mathfrak m\supseteq pA\)。商域 \(k=A/\mathfrak m\) 有限，因为 \(A/pA\) 由有限多个元素作为 \(\mathbb F_p\) 向量空间生成，而 \(k\) 是它的商。令
\[
D=\{\,\sigma\in G\mid\sigma(\mathfrak m)=\mathfrak m\,\},
\]
这是 \(\mathfrak m\) 在 \(G\) 作用下的稳定子群，故为子群。记商映射为 \(q:A\longrightarrow k=A/\mathfrak m\)。对 \(\sigma\in D\) 定义
\[
\rho(\sigma)\bigl(q(a)\bigr)=q\bigl(\sigma(a)\bigr)\qquad(a\in A).
\]
由于 \(\sigma(\mathfrak m)=\mathfrak m\)，此式与代表元无关，且逆映射由 \(\sigma^{-1}\) 给出；\(\sigma\) 固定整数，所以 \(\rho(\sigma)\) 固定 \(\mathbb F_p\)。于是得到同态 \(\rho:D\longrightarrow\operatorname{Gal}(k/\mathbb F_p)\)，记其像为 \(H\)。现在证明 \(H\) 是全群。

任取 \(x\in k^H\)。\(G\) 作用下 \(\mathfrak m\) 的轨道由有限个不同的极大理想组成，它们两两互素。\cref{b1:thm:crt}给出 \(a\in A\)，在 \(\mathfrak m\) 处的余数为 \(x\)，在其余轨道理想处的余数为零。多项式
\[
P(T)=\prod_{\sigma\in G}\bigl(T-\sigma(a)\bigr)
\]
的系数属于 \(A^G=\mathbb Z\)。若 \(\sigma\in D\)，因 \(x\) 被 \(H\) 固定，\(\sigma(a)\) 模 \(\mathfrak m\) 为 \(x\)；若 \(\sigma\notin D\)，则 \(a\) 模 \(\sigma^{-1}(\mathfrak m)\) 为零，故 \(\sigma(a)\) 模 \(\mathfrak m\) 为零。写 \(|D|=p^e u\)、\(p\nmid u\)，约化后得到
\[
\bar P(T)=T^{|G|-|D|}(T-x)^{|D|}
=T^{|G|-|D|}\left(T^{p^e}-x^{p^e}\right)^u.
\]
若直接看 \(T^{|G|-1}\) 的系数，只得到 \(-|D|x\)，而 \(p\) 可能整除 \(|D|\)。写成上式后，降次 \(p^e\) 的项保留下可逆的因子 \(u\)：\(T^{|G|-p^e}\) 的系数为 \(-u x^{p^e}\)。\(\bar P\) 的全部系数在 \(\mathbb F_p\) 中，故 \(x^{p^e}=c\in\mathbb F_p\)。于是 \((x-c)^{p^e}=0\)，得到 \(x=c\)。所以 \(k^H=\mathbb F_p\)。有限域扩张 \(k/\mathbb F_p\) 为 Galois 扩张，\cref{b1:thm:finite-galois-correspondence}迫使 \(H=\operatorname{Gal}(k/\mathbb F_p)\)。

因此有限域 \(k\) 的 Frobenius 自同构 \(\operatorname{Frob}_p:b\mapsto b^p\) 在 \(D\) 中有一个提升 \(\sigma\)。提升满足
\[
q\circ\sigma|_A=\operatorname{Frob}_p\circ q,
\qquad\text{即}\qquad
q\bigl(\sigma(\alpha_i)\bigr)=q(\alpha_i)^p\quad(1\leq i\leq n).
\]
等式 \(f(X)=\prod_i(X-\alpha_i)\) 经 \(q\) 约化后表明，\(q(\alpha_i)\) 是 \(\bar f\) 的全部根；因 \(\bar f\) 无重根，原根集到这些约化根的映射是双射。由\cref{b1:thm:finite-field-subfields}及\cref{b1:thm:finite-field-automorphisms}，次数为 \(d_j\) 的不可约因子的根恰组成长度 \(d_j\) 的 Frobenius 轨道：一个根生成的域有 \(p^{d_j}\) 个元素，而固定该根的最小正 Frobenius 次幂正是 \(d_j\)。上述相容等式使 \(\sigma\) 在原根集上的置换与 Frobenius 在约化根集上的置换对应，故循环长度也相同。
\end{proof}

因此，每次无重根的模素数分解都给出群中的一种循环型。将不同素数提供的信息与传递性、判别式判据结合，往往能确定整个群。\cref{b1:prop:quintic-x5-x-1}将用这一方法计算 \(X^5-X-1\) 的 Galois 群。

\ProseParagraphBreak
以先前的 \(f=X^3-2\) 为例，模 \(5\) 有
\[
\bar f=(X-3)(X^2+3X+4).
\]
二次因子的判别式为 \(3\in\mathbb F_5\)，不是平方，故因子次数为 \(1,2\)，且约化无重根。上述引理给出原 Galois 群中的一个换位；这与前面直接算出的 \(S_3\) 相符。这里算因子次数是在有限域中完成的，换位属于特征零的分裂域自同构群则依赖刚才的提升论证。

% !TeX root = ../../Book1.tex
% 纲要标题：### 19.4 中间域的合成、交与嵌入位置
% 纲要位置：计划纲要.md，第 666 行。

\section{中间域的合成、交与嵌入位置}
\label{b1:sec:1904}

% 纲要内容（仅作写作计划，尚非教材正文）：
%
% * 合成域与交域；解释一侧 Galois 假设、限制映射及反向的群运算，说明非正规扩张的次数公式可能失败
% * 同构类型与嵌入位置的区别；展开同构延拓与共轭子群的对应
% * 在正文证明正规化子商群给出中间域自同构群，并解释嵌入数、共轭子域数与自同构数的关系；相应习题改为四次分裂域的具体计算
%
% ---
%

\cref{b1:prop:cubic-splitting-field-correspondence}已将一个分裂域的中间域全部列出。现在考察两个中间域的合成与交，以及同构的域在同一个环境域中为何仍可能占据不同位置。

\subsection{合成域与交域}
\label{b1:subsec:1904:02}

在一个共同环境域中，两个子域 \(L,E\) 的合成域 \(LE\) 指包含 \(L\cup E\) 的最小子域。若 \(L/K\) 是一个可分多项式的分裂域，换基域到 \(E\) 后，把同一批根加入 \(E\) 得到的就是 \(LE\)。因此，要使 \(LE/E\) 成为 Galois 扩张，只需在 \(L/K\) 这一侧具有正规性与可分性，并不要求 \(E/K\) 也 Galois。

\begin{proposition}\label{b1:prop:galois-compositum}
在 \(K\) 的一个共同代数闭包内，设 \(L/K\) 是有限 Galois 扩张，\(E/K\) 是有限扩张，令 \(LE\) 为由 \(L\cup E\) 生成的子域，并记 \(D=L\cap E\)。则 \(LE/E\) 为 Galois 扩张，限制映射给出同构
\[
\operatorname{Gal}(LE/E)\cong\operatorname{Gal}(L/D).
\]
因此
\[
[LE:E]=[L:D],\qquad
[LE:K]=\frac{[L:K][E:K]}{[D:K]}.
\]
\end{proposition}
\begin{proof}
把 \(L\) 写为某个可分 \(K\) 多项式 \(f\) 的分裂域。把 \(f\) 看作 \(E[x]\) 中的多项式，它仍没有重根；添加其全部根到 \(E\) 后得到的正是 \(LE\)，所以 \(LE/E\) 为有限 Galois 扩张。若 \(\sigma\in\operatorname{Gal}(LE/E)\)，则它固定 \(E\)，因而也固定 \(K\subseteq E\)。其限制 \(\sigma|_L:L\rightarrow LE\) 是一个 \(K\) 嵌入，\(L/K\) 的正规性给出 \(\sigma(L)=L\)，所以才有同态
\[
\operatorname{res}:\operatorname{Gal}(LE/E)\longrightarrow\operatorname{Gal}(L/K),
\qquad \sigma\longmapsto\sigma|_L.
\]
若限制是恒等，\(\sigma\) 就逐点固定 \(L\cup E\)。由于 \(LE\) 由这两个域生成，它也逐点固定整个 \(LE\)，故 \(\sigma=1\)。因此限制映射单射，记像为 \(H\leq\operatorname{Gal}(L/K)\)。

对 \(a\in L\)，因为 \(H\) 是限制映射的像，它被 \(H\) 固定当且仅当它被 \(\operatorname{Gal}(LE/E)\) 的全部元素固定。这个群在 \(LE\) 中的不动域为 \(E\)，所以
\[
L^H=L\cap(LE)^{\operatorname{Gal}(LE/E)}=L\cap E=D.
\]
在 \(L/K\) 的 Galois 对应中，这迫使 \(H=\operatorname{Gal}(L/D)\)，因而限制映射满到所述子群。两边都是 Galois 扩张的群，取阶便有 \([LE:E]=[L:D]\)。再用塔式公式，
\[
[LE:K]=[LE:E][E:K]
=\frac{[L:K]}{[D:K]}[E:K],
\]
得到第二项次数等式。
\end{proof}
若 \(E_1,E_2\) 都是一个有限 Galois 扩张 \(M/K\) 的中间域，对应子群 \(H_1,H_2\)，则
\[
\operatorname{Gal}(M/E_1E_2)=H_1\cap H_2,\qquad
\operatorname{Gal}(M/E_1\cap E_2)=\langle H_1,H_2\rangle.
\]
第一式是因为逐点固定由 \(E_1\cup E_2\) 生成的较大域，等价于同时逐点固定 \(E_1\) 与 \(E_2\)，所以要取两个固定群的交。另一方面，被生成子群 \(\langle H_1,H_2\rangle\) 固定，等价于同时被 \(H_1\) 与 \(H_2\) 固定，因此
\[
M^{\langle H_1,H_2\rangle}=M^{H_1}\cap M^{H_2}=E_1\cap E_2.
\]
再用 Galois 对应的互逆性，得到第二式。合成域是包含两域的最小中间域，交域是包含于两域的最大中间域；反转包含关系后，它们分别对应子群的交与包含两群的最小子群。

\ProseParagraphBreak
取\cref{b1:fig:cubic-intermediate-fields}中的 \(E_1=\mathbb Q(a)\)、\(E_2=\mathbb Q(\zeta a)\)，并沿用\cref{b1:prop:cubic-splitting-field-correspondence}的记号 \(L,G,s,t\)。它们对应 \(H_1=\langle t\rangle\)、\(H_2=\langle s^2t\rangle\)。这两个二阶子群的交为 \(1\)，且两个不同的换位生成 \(G\cong S_3\)，故上述公式给出
\[
E_1E_2=L,\qquad E_1\cap E_2=\mathbb Q.
\]
前一等式也可由 \(a\) 与 \(\zeta a\) 的比得到 \(\zeta\) 而直接看出。因此 \([E_1E_2:\mathbb Q]=6\)，并非 \([E_1:\mathbb Q][E_2:\mathbb Q]=9\)：两个域都不正规时，交域为基域不能保证次数相乘。此处
\[
H_1H_2=\{1,t,s^2t,ts^2t\}=\{1,t,s^2t,s\}
\]
只有四个元素，依 Lagrange 定理不是六阶群 \(G\) 的子群；而 \(\langle H_1,H_2\rangle\) 还包含任意长的有限乘积，恰为 \(G\)。所以交域公式不能把生成子群替换为乘积集合。若其中一群正规，则在取积或取逆时，可以用正规性把另一群的元素移过它的元素，所得仍在乘积集合中，故乘积集合是子群。此时它包含两群，且每个元素都是两群元素的乘积，所以等于生成子群。

\subsection{同构类型与嵌入位置}
\label{b1:subsec:1904:03}

对两个实际子域 \(E_1,E_2\subseteq L\)，要求 \(E_1=E_2\) 是比较其中的元素；要求 \(E_1\cong_K E_2\) 则只要求存在固定 \(K\) 的域同构，允许它们在 \(L\) 中的位置不同。例如，\(\mathbb Q(a)\) 和 \(\mathbb Q(\zeta a)\) 都同构于 \(\mathbb Q[x]/(x^3-2)\)，存在保持 \(\mathbb Q\) 并把 \(a\) 送到 \(\zeta a\) 的同构，但它们是分裂域中的不同子域。

\ProseParagraphBreak
在有限 Galois 扩张 \(L/K\) 中，记 \(G=\operatorname{Gal}(L/K)\)。还有第三个问题：是否存在 \(\sigma\in G\) 使 \(\sigma(E_1)=E_2\)？若有，限制 \(\sigma|_{E_1}:E_1\rightarrow E_2\) 就是 \(K\) 同构。反过来，给定 \(K\) 同构 \(\varphi:E_1\overset{\sim}{\longrightarrow}E_2\)，把 \(E_2\) 包含到一个含 \(L\) 的代数闭包 \(\Omega\) 中，便可将 \(\varphi\) 看作 \(K\) 嵌入 \(E_1\rightarrow\Omega\)。由\cref{b1:thm:embedding-extension}，它能延拓为 \(\widetilde\varphi:L\rightarrow\Omega\)。\(L/K\) 的正规性给出 \(\widetilde\varphi(L)=L\)，所以 \(\widetilde\varphi\in G\)，且 \(\widetilde\varphi(E_1)=E_2\)。因此在这个有限 Galois 环境中，\(K\) 同构与被 \(G\) 中元素互相送到等价，子域相等则是更强的要求。

\ProseParagraphBreak
若 \(H_i=\operatorname{Gal}(L/E_i)\)，由 \(\sigma(E_1)=E_2\) 及\cref{b1:thm:galois-quotient}证明中的共轭公式，
\[
E_2=\sigma(L^{H_1})=L^{\sigma H_1\sigma^{-1}},\qquad
H_2=\sigma H_1\sigma^{-1}.
\]
反过来，共轭子群的共同不动域也由 \(\sigma\) 互相送到。Galois 对应枚举的是实际子域，而按 \(K\) 同构类型分类则须再把共轭子群归为一类。

\begin{proposition}\label{b1:prop:subfield-normalizer-automorphisms}
设 \(L/K\) 为有限 Galois 扩张，\(G=\operatorname{Gal}(L/K)\)，\(H\leq G\)，\(E=L^H\)，并记
\[
N_G(H)=\{\,\sigma\in G\mid\sigma H\sigma^{-1}=H\,\}.
\]
则限制映射给出
\[
\operatorname{Aut}_K(E)\cong N_G(H)/H.
\]
域 \(L\) 中与 \(E\) 在 \(K\) 上同构的不同中间域共有 \([G:N_G(H)]\) 个，而且
\[
[E:K]=[G:N_G(H)]\,|\operatorname{Aut}_K(E)|.
\]
\end{proposition}
\begin{proof}
由\cref{b1:thm:galois-quotient}证明中的公式 \(\sigma(E)=L^{\sigma H\sigma^{-1}}\) 及 Galois 对应的单射性，\(\sigma(E)=E\) 当且仅当 \(\sigma\in N_G(H)\)。因此限制给出群同态
\[
\operatorname{res}:N_G(H)\longrightarrow\operatorname{Aut}_K(E),\qquad
\sigma\longmapsto\sigma|_E,
\]
核为逐点固定 \(E\) 的群 \(\operatorname{Gal}(L/E)=H\)；由正规化子的定义，\(H\) 在 \(N_G(H)\) 中正规。任取 \(\tau\in\operatorname{Aut}_K(E)\)，将它看作到一个含 \(L\) 的代数闭包 \(\Omega\) 的 \(K\) 嵌入，由\cref{b1:thm:embedding-extension}延拓为 \(\widetilde\tau:L\rightarrow\Omega\)。由于 \(L/K\) 正规，\(\widetilde\tau(L)=L\)，故 \(\widetilde\tau\in G\)。它在 \(E\) 上等于 \(\tau\)，所以 \(\widetilde\tau(E)=E\)，即属于 \(N_G(H)\)。限制映射因而满射，第一同构定理给出所述商群同构。

\(G\) 作用在 \(L\) 的中间域集合上，作用为 \(E'\mapsto\sigma(E')\)。由前面的同构延拓论证，与 \(E\) 在 \(K\) 上同构的域恰为这个作用下 \(E\) 的轨道；将 \(E\) 保持为集合的稳定子群恰为 \(N_G(H)\)。\cref{b1:thm:orbit-stabilizer}给出不同域的数目 \([G:N_G(H)]\)。最后，由 Galois 对应与指数乘法公式，
\[
[E:K]=[G:H]=[G:N_G(H)]\,[N_G(H):H]
=[G:N_G(H)]\,|\operatorname{Aut}_K(E)|.
\]
\end{proof}
这个次数等式也可以直接按嵌入来理解。中间扩张 \(E/K\) 可分，所以它到代数闭包的不同 \(K\) 嵌入共有 \([E:K]\) 个。每个嵌入都延拓到 \(L\)，由正规性，其像是 \(L\) 内与 \(E\) 同构的某个子域。固定其中一个像域 \(E'\) 并选一个 \(K\) 同构 \(\varphi:E\rightarrow E'\) 后，全部以 \(E'\) 为像的嵌入恰为 \(\varphi\circ\tau\)，其中 \(\tau\in\operatorname{Aut}_K(E)\)。因此嵌入总数等于共轭副本的数目乘以每个副本内的同构数。自同构要求像仍为原来的 \(E\)，一般嵌入则可以把它送到另一个副本。

\ProseParagraphBreak
在\cref{b1:prop:cubic-splitting-field-correspondence}的三次分裂域中，对 \(E=\mathbb Q(a)\)，对应 \(H=\langle t\rangle\)。正规化 \(H\) 的元素须在共轭下固定它唯一的非单位元 \(t\)；在 \(S_3\) 中，与这个换位交换的元素只有 \(1,t\)，所以 \(N_G(H)=H\)。于是 \(E\) 有三个不同的共轭子域，却只有恒等 \(\mathbb Q\) 自同构，次数公式为 \(3=3\cdot1\)。对 \(E=\mathbb Q(\zeta)\)，对应的三阶子群 \(H=\langle s\rangle\) 正规，故 \(N_G(H)=G\)。这个二次域是同一环境中唯一的同构副本，但有两个 \(\mathbb Q\) 自同构，次数公式为 \(2=1\cdot2\)。两种情形都具有次数那么多的嵌入，只是嵌入是否保持同一子域不同。


\begin{chaptersummary}
有限自同构群的不动域总给出一个次数恰为群阶的 Galois 扩张。由此，中间域与子群反向对应，扩张次数对应群指数，对基域为 Galois 的中间域对应正规子群，限制映射给出相应商群。具体计算时，群的抽象类型还不够：须明确它作用于哪些根、哪些元素被固定。合成域与交域公式需要核对正规可分假设，非正规扩张中的交域为基域并不保证次数相乘。同构的中间域也可能在同一个分裂域中占据不同位置；若 \(E=L^H\)，整体保持 \(E\) 的群为 \(N_G(H)\)，逐点固定 \(E\) 的群为 \(H\)，所以 \(\operatorname{Aut}_K(E)\cong N_G(H)/H\)。下一章将研究 Galois 群为循环群的扩张，考察生成自同构、范数条件与扩张生成元之间的关系。
\end{chaptersummary}

\BookExercises

\begin{exercise}\label{b1:ex:19-characteristic-two}
比较 \(\mathbb F_4/\mathbb F_2\) 与 \(\mathbb F_2(u)/\mathbb F_2(u^2)\)，其中 \(u\) 超越。它们同为二次扩张，分别是否为 Galois 扩张？求各自的基域自同构群。
\end{exercise}
\begin{solution}
前者是 \(x^2+x+1\) 的分裂域，该多项式导数为一，故可分。Galois 群由恒等与平方映射组成，为 \(C_2\)。后者为纯不可分扩张；若嵌入固定 \(u^2\)，则 \(u\) 的像必须满足 \(x^2-u^2=(x-u)^2=0\)，只能仍为 \(u\)。所以自同构群平凡，扩张不可分，不是 Galois 扩张。这个最小多项式虽有重根，却已在扩张域中完全分裂，且其根 \(u\) 生成扩张，所以该扩张仍正规。二次代数扩张的正规性，并不保证它在特征二时也可分。
\end{solution}

\begin{exercise}\label{b1:ex:19-real-fourth-root}
令 \(a=\sqrt[4]2>0\)、\(E=\mathbb Q(a)\)。求 \(\operatorname{Aut}_{\mathbb Q}(E)\) 及其不动域，解释为何这个不动域不是 \(\mathbb Q\)。
\end{exercise}
\begin{solution}
\(x^4-2\) 由素数二的 Eisenstein 判据不可约，是 \(a\) 的最小多项式，故 \([E:\mathbb Q]=4\)。\(E\subseteq\mathbb R\)，所以 \(a\) 的像只能是两个实根 \(a,-a\)。将 \(\mathbb Q[x]/(x^4-2)\) 的生成元分别送到这两个根，所得的像都为 \(E\)，所以确有这两个自同构。因此 \(H=\operatorname{Aut}_{\mathbb Q}(E)\cong C_2\)，其非平凡元 \(s\) 把 \(a\) 变号。

\(\mathbb Q(a^2)=\mathbb Q(\sqrt2)\) 被 \(H\) 固定，且 \([E:\mathbb Q(a^2)]=4/2=2\)。\cref{b1:thm:artin-fixed-field}给出 \([E:E^H]=|H|=2\)，由包含关系和塔式公式可知 \(E^H=\mathbb Q(\sqrt2)\)。该定理保证的是 \(E/\mathbb Q(\sqrt2)\) 为 Galois 扩张；原来的 \(E/\mathbb Q\) 缺少非实共轭根，仍不正规。这也具体说明，在非 Galois 扩张中，\(\mathbb Q\subseteq E^{\operatorname{Aut}_{\mathbb Q}(E)}\) 可以是严格包含。
\end{solution}

\begin{exercise}\label{b1:ex:19-rational-inversion}
设 \(k\) 为任意域，\(t\) 在 \(k\) 上超越。令 \(s\in\operatorname{Aut}_k(k(t))\) 为满足 \(s(t)=t^{-1}\) 的 \(k\) 自同构。求 \(\langle s\rangle\) 的不动域。
\end{exercise}
\begin{solution}
\(s\) 逐点固定 \(k\)，且 \(s^2(t)=t\)，故 \(s^2=\operatorname{id}_{k(t)}\)。又 \(s\ne\operatorname{id}\)，因为 \(t=t^{-1}\) 将使超越元满足 \(t^2=1\)，所以 \(\langle s\rangle\) 的阶为二。令 \(v=t+t^{-1}\)，它被 \(s\) 固定，且 \(t\) 满足 \(x^2-vx+1=0\)，所以 \([k(t):k(v)]\leq2\)。若次数为一，\(s\) 固定 \(v\) 与 \(k\) 就固定 \(t\)，矛盾，故次数为二。与\cref{b1:thm:artin-fixed-field}的次数比较，得到不动域恰为 \(k(t+t^{-1})\)。论证在特征二也成立。
\end{solution}

\begin{exercise}\label{b1:ex:19-faithful-action}
让抽象循环群 \(C_4\) 的生成元通过复共轭作用在 \(\mathbb C\) 上。求不动域，并解释为何不能据此说 \([\mathbb C:\mathbb R]=4\)。
\end{exercise}
\begin{solution}
共同不动元为 \(\mathbb R\)。但作用同态 \(\rho:C_4\rightarrow\operatorname{Aut}(\mathbb C)\) 的核由生成元的平方生成，其像只有恒等和共轭两个元素。正如群作用通过核的商群给出忠实作用，这里的不动元只取决于像群 \(\rho(C_4)\cong C_4/\ker\rho\cong C_2\)。不动域定理计算的是实际互异自同构所组成的群的阶，因而给出 \([\mathbb C:\mathbb R]=2\)。只有当抽象群的作用忠实时，才能直接把它的阶作为这个次数。
\end{solution}

\begin{exercise}\label{b1:ex:19-biquadratic-five}
列出 \(L=\mathbb Q(\sqrt2,\sqrt5)\) 的全部中间域，并写明各自对应的子群。
\end{exercise}
\begin{solution}
若 \(\sqrt5=a+b\sqrt2\)，平方比较给出 \(ab=0\) 及 \(a^2+2b^2=5\)，无有理解，因为五与 \(5/2\) 均非有理数平方。故总次数为四，群由分别变号的 \(s,t\) 生成，满足 \(s^2=t^2=1,st=ts\)。子群为平凡群、全群和 \(\langle s\rangle,\langle t\rangle,\langle st\rangle\)。若 \(s\) 改变 \(\sqrt2\) 的符号、\(t\) 改变 \(\sqrt5\) 的符号，则对应域依次为 \(L,\mathbb Q,\mathbb Q(\sqrt5),\mathbb Q(\sqrt2),\mathbb Q(\sqrt{10})\)。这些子群都是正规子群，所以各中间域对 \(\mathbb Q\) 均为 Galois 扩张。
\end{solution}

\begin{exercise}\label{b1:ex:19-cyclic-twelve}
设有限 Galois 扩张 \(L/K\) 的群为 \(C_{12}=\langle\sigma\rangle\)。求所有可能的中间域次数，判断每个次数的中间域是否唯一，并求四次中间域对 \(K\) 的 Galois 群。
\end{exercise}
\begin{solution}
循环群对每个正因子有唯一相应阶的子群，所以中间域次数为 \(1,2,3,4,6,12\)，每个次数恰有一个。设四次中间域对应 \(H\)，则 \([G:H]=4\)，而 \(|H|=12/4=3\)：对应的是指数为四、阶为三的子群，不是四阶子群。因此 \(H=\langle\sigma^4\rangle\)，且它正规。由\cref{b1:thm:galois-quotient}，四次域的 Galois 群为 \(C_{12}/\langle\sigma^4\rangle\cong C_4\)。
\end{solution}

\begin{exercise}\label{b1:ex:19-mod-primes-s4}
设 \(f=X^4-X-1\in\mathbb Q[X]\)。求 \(f\) 在 \(\mathbb Q\) 上的 Galois 群，并证明所得结论。
\end{exercise}
\begin{solution}
在 \(\mathbb F_2[X]\) 中，\(\bar f=X^4+X+1\) 无根，且除以唯一的首一不可约二次式 \(X^2+X+1\) 余 \(1\)，故 \(\bar f\) 不可约，因而无重根。由于 \(f\) 首一，模 \(2\) 不可约也推出 \(f\) 在 \(\mathbb Q[X]\) 中不可约。循环型判据给出 Galois 群 \(G\leq S_4\) 中一个四轮换。

在 \(\mathbb F_7[X]\) 中，
\[
\bar f=(X-3)(X^3+3X^2+2X-2).
\]
右侧三次因子在 \(0,1,\ldots,6\) 处的值依次为 \(5,4,1,2,6,5,5\)，故在 \(\mathbb F_7\) 中无根，从而不可约；它在 \(3\) 处的值为 \(2\ne0\)，所以约化无重根，判据又给出一个三轮换。由 Lagrange 定理，四轮换和三轮换分别给出 \(4\bigm\vert|G|\) 与 \(3\bigm\vert|G|\)。由于三、四互素，故 \(12\bigm\vert|G|\)；而 \(G\leq S_4\) 又给出 \(|G|\bigm\vert24\)。因此 \(|G|\) 只能为十二或二十四，这并没有断言那两个元素生成的子群恰有十二个元素。

若 \(|G|=12\)，四轮换为奇置换，所以符号同态 \(\operatorname{sgn}|_G:G\rightarrow\{\pm1\}\) 满射，核为 \(G\cap A_4\)。第一同构定理给出 \(|G\cap A_4|=|G|/2=6\)，与\cref{b1:prop:alternating-three-cycles}中 \(A_4\) 没有六阶子群的结论矛盾。因此 \(|G|=24\)，即 \(G=S_4\)。
\end{solution}

\begin{exercise}\label{b1:ex:19-dihedral-quadratics}
在 \(L=\mathbb Q(a,i)\)、\(a^4=2\) 中，设 \(r(a)=ia,r(i)=i\)，\(s\) 为复共轭。证明三个二次中间域为 \(\mathbb Q(i),\mathbb Q(\sqrt2),\mathbb Q(\sqrt{-2})\)，并写出对应子群。
\end{exercise}
\begin{solution}
二次域对应 \(D_8\) 中的指数二子群，即四阶子群。它们恰为 \(\langle r\rangle\)、\(\langle r^2,s\rangle\)、\(\langle r^2,rs\rangle\)：若四阶子群含奇次幂的 \(r\)，就包含 \(\langle r\rangle\)；否则其旋转部分为 \(\{1,r^2\}\)，两个反射的指数奇偶相同，只给出后两种。三群分别固定 \(i,a^2,ia^2\)，这些元素生成的域均为二次域；次数相同使包含成为相等，因此得所列三个域。
\end{solution}

\begin{exercise}\label{b1:ex:19-cyclotomic-five}
令 \(\zeta\) 为本原五次单位根。求 \(\mathbb Q(\zeta)\) 的唯一二次中间域，并写出其一个生成元满足的二次方程。
\end{exercise}
\begin{solution}
正文已证明群为 \(C_4\)，故二次中间域唯一，对应由复共轭生成的二元子群。令 \(u=\zeta+\zeta^{-1}\)，则其为实数且被共轭固定。由 \(1+\zeta+\zeta^2+\zeta^3+\zeta^4=0\)，有 \(\zeta^2+\zeta^{-2}=-1-u\)，故 \(u^2=1-u\)。多项式 \(t^2+t-1\) 无有理根，因而 \(\mathbb Q(u)\) 为二次域，且 \((2u+1)^2=5\)。所求域为 \(\mathbb Q(\sqrt5)\)。
\end{solution}

\begin{exercise}\label{b1:ex:19-reducible-permutation}
求 \((x^2-2)(x^2-3)\) 对 \(\mathbb Q\) 的 Galois 群及其在四个根上的轨道。将它与不可约四次多项式 \(x^4-10x^2+1\) 的情形比较。
\end{exercise}
\begin{solution}
二者分裂域都是 \(\mathbb Q(\sqrt2,\sqrt3)\)，抽象 Galois 群均为 \(C_2\times C_2\)。前一个多项式的根分成两个轨道 \(\{\sqrt2,-\sqrt2\}\)、\(\{\sqrt3,-\sqrt3\}\)，因为有理系数自同构不能把平方为二的元素送到平方为三的元素。后一个多项式的四根为 \(\pm\sqrt2\pm\sqrt3\)，独立变号使其成为一个轨道。这说明相同的抽象群可以在不同根集上给出不同的置换作用。
\end{solution}

\begin{exercise}\label{b1:ex:19-cyclic-cubic}
证明 \(f=x^3-3x-1\) 在 \(\mathbb Q\) 上的 Galois 群为 \(C_3\)。计算判别式时可用全部根 \(a_i\) 上的导数乘积。
\end{exercise}
\begin{solution}
有理根只能为 \(\pm1\)，二者都不是根，故三次式不可约。设三个根为 \(a_1,a_2,a_3\)。由 \(f'(a_i)=\prod_{j\ne i}(a_i-a_j)\)，对三个指标相乘得到 \(\prod_i f'(a_i)=-\operatorname{disc}(f)\)。又 \(f'(x)=3(x-1)(x+1)\)，所以
\[
\prod_i f'(a_i)=27\Bigl(\prod_i(a_i-1)\Bigr)\Bigl(\prod_i(a_i+1)\Bigr)
=27\bigl(-f(1)\bigr)\bigl(-f(-1)\bigr)=-81.
\]
故判别式为 \(81\)，是非零平方。\eqref{b1:eq:discriminant-alternating-group}说明群包含于 \(A_3\)，而不可约性使作用传递，群阶被三整除。因此它就是 \(A_3\cong C_3\)。
\end{solution}

\begin{exercise}\label{b1:ex:19-product-groups}
设 \(L/K,E/K\) 都是有限 Galois 扩张，且 \(L\cap E=K\)。证明 \(LE/K\) 的 Galois 群通过限制同构于 \(\operatorname{Gal}(L/K)\times\operatorname{Gal}(E/K)\)。
\end{exercise}
\begin{solution}
两域都是可分多项式的分裂域。取这两个多项式的全部互异不可约因子的乘积，它仍可分，且分裂域为 \(LE\)，故 \(LE/K\) 为 Galois 扩张。两个限制给出到直积的同态，固定两域的自同构固定其合成域，所以该同态单射。由\cref{b1:prop:galois-compositum}以及交域假设，\([LE:K]=[L:K][E:K]\)。这恰为目标直积群的阶，也等于来源群的阶，所以单射为满射。
\end{solution}

\begin{exercise}\label{b1:ex:19-two-fourth-root-fields}
设 \(a=\sqrt[4]2>0\)，比较 \(E_1=\mathbb Q(a)\) 与 \(E_2=\mathbb Q(ia)\)。证明二者在 \(\mathbb Q\) 上同构而不相等，并求它们的交域、合成域。
\end{exercise}
\begin{solution}
\(a,ia\) 都以不可约多项式 \(x^4-2\) 为最小多项式，所以代入映射给出两个域与 \(\mathbb Q[x]/(x^4-2)\) 的同构。但前者为实域，后者含非实元，故不相等。两者均含 \(a^2=\sqrt2\)，因为 \((ia)^2=-a^2\)。交域对 \(\mathbb Q\) 的次数整除四且至少为二；若为四，则两域相等，矛盾。因此交域恰为 \(\mathbb Q(\sqrt2)\)。合成域含 \((ia)/a=i\)，故为 \(\mathbb Q(a,i)\)。由于 \(\mathbb Q(a)\) 是实域，\(i\notin\mathbb Q(a)\)，而 \(i\) 满足 \(x^2+1=0\)，所以合成域对 \(\mathbb Q(a)\) 的次数为二，总次数为八。

这里直接识别了合成域并独立计算其次数。\(E_1/\mathbb Q\) 与 \(E_2/\mathbb Q\) 都不正规，不能以\cref{b1:prop:galois-compositum}为依据套用次数乘除公式；本题数值恰与 \(4\cdot4/2\) 相同，并不使缺少假设的公式成立。
\end{solution}

\begin{exercise}\label{b1:ex:19-normalizer-automorphisms}
设 \(a=\sqrt[4]2>0\)、\(L=\mathbb Q(a,i)\)，并令 \(r(a)=ia,r(i)=i\)，\(s\) 为复共轭。由正文知 \(G=\operatorname{Gal}(L/\mathbb Q)=\langle r,s\rangle\cong D_8\)。对 \(E=\mathbb Q(a)\)，求逐点固定 \(E\) 的子群 \(H\) 及其正规化子，写出 \(E\) 的全部 \(\mathbb Q\) 自同构和 \(L\) 中与 \(E\) 在 \(\mathbb Q\) 上同构的全部中间域，并核对
\[
[E:\mathbb Q]=[G:N_G(H)]\,|\operatorname{Aut}_{\mathbb Q}(E)|.
\]
再对 \(E'=\mathbb Q(i)\) 完成相同的计算，并比较两个子域的嵌入如何分成不同像域。
\end{exercise}
\begin{solution}
\(s\) 固定 \(a\)，故 \(\langle s\rangle\subseteq H\)。又 \([L:E]=2\)，Galois 对应给出 \(|H|=2\)，所以 \(H=\langle s\rangle\)。由于 \(H\) 只有一个非单位元，将它保持为集合就须固定这个元素。利用 \(srs=r^{-1}\)，有
\[
r^j s r^{-j}=r^{2j}s\qquad(0\leq j<4).
\]
元素 \(r^js\) 的共轭作用在 \(s\) 上也给出 \(r^{2j}s\)，因此当且仅当 \(j\) 为偶数时属于正规化子。于是
\[
N_G(H)=\{1,r^2,s,r^2s\}=\langle r^2,s\rangle.
\]
\cref{b1:prop:subfield-normalizer-automorphisms}给出 \(\operatorname{Aut}_{\mathbb Q}(E)\cong N_G(H)/H\cong C_2\)。两自同构分别把 \(a\) 送到 \(a\) 与 \(-a\)，非平凡者来自 \(r^2\) 的限制。

不同共轭子域共有 \([G:N_G(H)]=8/4=2\) 个，实际为 \(\mathbb Q(a)\) 和 \(\mathbb Q(ia)\)：\(G\) 将 \(a\) 送到 \(\pm a,\pm ia\)，正负号不改变生成的域。这两个域不同，因为前者为实域，后者含非实元。四个 \(\mathbb Q\) 嵌入因而分成两组，每组有两个嵌入，其像为同一个子域，次数公式在这里成为 \(4=2\cdot2\)。

对 \(E'=\mathbb Q(i)\)，\(r\) 固定 \(i\)，而 \([L:E']=4\)，故对应子群为 \(H'=\langle r\rangle\)。它在 \(G\) 中正规，因而 \(N_G(H')=G\)。相应自同构群为 \(G/H'\cong C_2\)，其元素分别将 \(i\) 送到 \(i\) 与 \(-i\)，共轭子域只有 \(E'\) 自身。此时
\[
[E':\mathbb Q]=2=1\cdot2.
\]
两个嵌入都保持同一子域，正是这个二次扩张正规性的表现。
\end{solution}

\begin{exercise}\label{b1:ex:19-translation-fixed-field}
设 \(k\) 为特征 \(p>0\) 的域，\(t\) 在 \(k\) 上超越。对每个 \(c\in\mathbb F_p\)，令 \(\sigma_c\in\operatorname{Aut}_k(k(t))\) 为逐点固定 \(k\)、满足 \(\sigma_c(t)=t+c\) 的自同构。求这些自同构所成群的不动域。
\end{exercise}
\begin{solution}
这 \(p\) 个自同构互异，且 \(\sigma_c\sigma_d=\sigma_{c+d}\)。元素 \(v=t^p-t\) 被全部自同构固定，因为 \(c^p=c\)。又 \(t\) 满足 \(x^p-x-v=0\)，所以 \([k(t):k(v)]\leq p\)。由\cref{b1:thm:artin-fixed-field}，对真正不动域的扩张次数恰为 \(p\)；由于 \(k(v)\) 包含于该不动域，塔式公式迫使两者相等。故不动域为 \(k(t^p-t)\)。
\end{solution}
